\chapter*{Preface}
\addcontentsline{toc}{chapter}{Preface}

\lettrine{T}{his booklet} is a self-contained excerpt of \BookletSubtitle
--- a book about arrays that starts from a voltage difference across a
transistor and goes up to tensor cores and quantum state vectors. Before
any of that is reachable, one question has to be answered first: what is
data, actually, and how does something abstract end up carried by
something physical?

That is the question this booklet follows, across three chapters. The
first asks what data is, using a short passage from Marx's \textit{Eighteenth
Brumaire} to show that the same fixed string of characters yields entirely
different data to a child, a historian, and a political economist ---
and works through the DIKW pyramid (Data, Information, Knowledge, Wisdom)
as one way of naming that gap, before arriving at a harder question: how
a word, a number, or a sign comes to stand in for something else at all.
The second takes up that question directly, and asks where representation
comes from in the first place, drawing on the hippocampus, episodic
memory, and four criteria for what counts as a representation, before
closing with a further question of its own: whether representation can
exist as a physical system, and not merely a mental one. The third picks
up exactly there, and follows that question to its end --- from an
abstract stand-in for something else, to a physical signal that can
actually carry it.

\vfill

\begin{flushright}
\textit{\BookletAuthor} \\
\textit{\BookletYear}
\end{flushright}
